O WebAssembly tem se consolidado como uma alternativa para execução de aplicações de alto desempenho em navegadores \emph{web}, permitindo a utilização de linguagens compiladas além da linguagem JavaScript. Entre elas, o Rust destaca-se pela eficiência computacional e sua segurança de memória, tornando-se uma opção importante para o desenvolvimento de aplicações executadas no ambiente \emph{web}. Neste contexto, este trabalho realiza uma análise comparativa entre a execução nativa de algoritmos implementados em Rust e a execução dos mesmos algoritmos compilados para WebAssembly. Para a realização dos experimentos, foram selecionados algoritmos representativos de diferentes classes de problemas computacionais a partir do The Computer Language Benchmarks Game. As implementações foram adaptadas para os dois ambientes de execução e avaliadas em ambiente controlado por contêineres Docker, com análise de tempo de execução e consumo de memória. Os resultados mostram que o Rust nativo foi mais rápido em todos os algoritmos avaliados, com \emph{overhead} do WebAssembly variando de 10,9\% a 126,9\% conforme o algoritmo, com mediana de 21,2\%. Em relação à memória, o ambiente WebAssembly manteve cerca de 210 MiB em repouso, devido ao navegador e à ferramenta de automação, o que tornou necessário o uso das métricas de \emph{baseline} e consumo líquido para isolar o consumo dos algoritmos, que foi reduzido na maioria dos casos e maior no mandelbrot, nos dois ambientes, em valor compatível com o \emph{buffer} de saída do algoritmo. Conclui-se que o WebAssembly pode ser uma alternativa viável para aplicações Rust executadas no navegador, mas que o custo adicional depende fortemente da carga de trabalho, devendo ser avaliado para cada aplicação quando o desempenho for um requisito crítico. Os resultados contribuem para a compreensão dos impactos do WebAssembly em aplicações Rust e auxiliam na tomada de decisões sobre quando utilizar essa tecnologia em aplicações executadas no navegador. 

\textbf{Palavras-chave:} Rust; WebAssembly; desempenho; benchmarking; execução nativa; Docker; Puppeteer.
